\section{Small examples and useful background}
\label{app:examples}

\subsection{A complete two-by-two system}

Take \(m=1,n=4\), with intervals \([0,4]\) and \([1,3]\).
Here \(d_4=12\), the interval lengths are \(4,2\), and
\[
C_n=1!\cdot12\cdot4!\cdot2!=576.
\]
The exact coefficient matrix is
\begingroup
\small
\setlength{\arraycolsep}{9pt}
\[
 A_{1,4}(x)=\left[\begin{matrix}4 x^{4} + 112 x^{3} - 112 x - 4 & - 48 x^{3} - 144 x^{2} - 48 x\\- 160 x^{3} - 144 x^{2} + 288 x + 16 & 48 x^{3} + 288 x^{2} + 144 x\end{matrix}\right] .
\]

\endgroup
Its determinant is
\[
\det A_{1,4}(x)=192x(x-1)^6.
\]
The exponent \(6=4+2\) is \(N_n\), while the factor \(x\) is
\(x^{(4-4)+(4-3)}\). The two remainders are
\[
R_{0,4}(x)=A_{0,0,4}(x)+A_{0,1,4}(x)\log x,\quad
R_{1,4}(x)=A_{1,0,4}(x)+A_{1,1,4}(x)\log x.
\]
Their first possible orders at \(x=1\) are respectively seven and
six, in agreement with the residue-at-infinity calculation.
At \(x=\ii\) the determinant is nonzero.

This example detects a normalization error immediately.
Using \(d_4/(4!2!)=1/4\) instead of \(576\) would divide the whole
matrix by \(2304\). Its \(x^4\) coefficient in the upper-left entry
would be \(1/576\), rather than the integer \(4\).

\subsection{Why the minimal polynomial cannot be skipped}

Consider \(y=\ii\) and \(U(Y)=Y^2+\ii Y-1\).
Then \(U(\ii)=-3\ne0\), but
\[
Y^3+\ii=(Y-\ii)(Y^2+\ii Y-1).
\]
The product of \(U\) over all three roots of \(Y^3+\ii\) is zero.
Nonvanishing at one root alone would not justify a nonzero
full-binomial norm. The minimal polynomial of \(y\) over
\(\QQ(\ii)\) is \(Y-\ii\); its resultant with \(U\) is \(-3\).
The same logical issue arises for perfect-power rational
approximations. Lemma~\ref{lem:resultant} handles it uniformly.

\subsection{Divisibility before and after evaluation}

A prime dividing every coefficient of a polynomial also divides its
value at \(\ii\) in \(\ZZ[\ii]\). The converse fails. Evaluation can
create cancellations, and an individual \(p\)-adic unit coefficient
need not remain visible in the sum.
For example, \(P(x)=1+x^2\) has unit coefficients at every odd prime,
but \(P(\ii)=0\). Less extreme cancellations can raise valuations
without making the value zero.
Thus boundary coefficient calculations alone cannot bound the gcd
of an evaluated row from above.

Our balanced arc minimum \(f=0\) likewise proves only that the
particular coefficientwise divisor has \(\gamma=0\). It is not a
proof that every possible evaluated-row divisor has rate zero.

\subsection{The least common multiple as an arithmetic cost}

Prime factorization gives
\[
\log d_n=\sum_{p^a\le n}\log p=n+o(n).
\]
The last equality uses the prime number theorem. Thus \(d_n^m\)
has exponential cost \(m\), although it is much smaller than
\((n!)^m\). The interval factorials in \(C_n\), by contrast,
cancel the factorial scale of the contour denominator.
Their position is crucial before Stirling's formula is applied.

Factors such as \(m!\) can be enormous but are fixed as \(n\)
grows. Their logarithms divided by \(n\) vanish.
For large \(m\), small-\(n\) regressions may confuse these constants
with exponential improvement. Exact identities avoid this ambiguity.

\subsection{A two-layer arc calculation by hand}

For \(L=14,z=3/40\), the outer zero arc is
\([0,\{14z\}]=[0,1/20]\), and the inner one is
\([\{z\},\{13z\}]=[3/40,39/40]\).
They are disjoint, as drawn in Figure~\ref{fig:arcs}.
With weights \(128,128\), at least one type contributes \(128\)
at every pole location. Formula~\eqref{eq:twogamma} collects
all such intervals and every periodic prime range.
